% !TeX root = ../informe.tex
% ===========================================================================
%  5. Modelo de datos
%
%  Lo que se evalúa aquí no es solo que las tablas estén normalizadas, sino
%  que el reparto de los datos entre las piezas del sistema esté razonado.
% ===========================================================================

\section{Modelo de datos}
\label{sec:modelo-datos}

\subsection{Estrategia de persistencia}
\label{sec:estrategia-datos}

\instruccion{Explica el reparto: qué almacén pertenece a qué pieza y qué
principio se sigue. Si se tomó un atajo respecto del diseño ideal — una sola
instancia del motor en lugar de varias, por ejemplo — dilo aquí y
justifícalo: es una pregunta previsible en la defensa.}

\ph{Cómo se reparten los datos, qué implica ese reparto y cómo se garantiza.}

\begin{table}[H]
  \centering
  \caption{Reparto de almacenes de datos.}
  \label{tab:bases}
  \begin{tabular}{L{3.2cm} L{3.4cm} L{6.6cm}}
    \toprule
    \textbf{Almacén} & \textbf{Pieza dueña} & \textbf{Contenido} \\
    \midrule
    \id{\ph{...}} & \id{\ph{...}} & \ph{...} \\
    \id{\ph{...}} & \id{\ph{...}} & \ph{...} \\
    \bottomrule
  \end{tabular}
\end{table}

\subsection{Diagrama entidad-relación}
\label{sec:der}

\instruccion{Un DER por almacén, o uno global con las fronteras marcadas. Lo
segundo comunica mejor la independencia de los datos, que es justo lo que se
evalúa. La figura se espera en \texttt{figuras/modelo-datos.pdf}.}

\begin{figure}[H]
  \centering
  \diagrama{modelo-datos}
  \caption{Modelo entidad-relación, con la frontera de cada almacén.}
  \label{fig:der}
\end{figure}

\subsection{Diccionario de datos}
\label{sec:diccionario}

\instruccion{Una tabla por entidad. No hace falta repetir todos los tipos
SQL: basta el campo, su tipo lógico y qué significa. Los tipos exactos están
en los scripts del anexo. Repite la tabla una vez por entidad relevante.}

\begin{table}[H]
  \centering
  \caption{Entidad \id{\ph{nombre}}.}
  \label{tab:dic-entidad}
  \begin{tabular}{L{3.4cm} L{3.2cm} L{6.6cm}}
    \toprule
    \textbf{Campo} & \textbf{Tipo} & \textbf{Descripción} \\
    \midrule
    \ph{...} & \ph{...} & \ph{...} \\
    \bottomrule
  \end{tabular}
\end{table}

\subsection{Integridad entre almacenes}
\label{sec:integridad}

\instruccion{Es el apartado que más distingue un diseño distribuido de un
monolito con varias tablas. Explica qué referencias cruzan la frontera entre
almacenes, por qué no pueden ser claves foráneas y dónde se valida entonces
la integridad. Si el sistema tiene un solo almacén, sustituye este apartado
por las restricciones de integridad referencial y dilo.}

\ph{Referencias lógicas sin clave foránea: cuáles son, dónde se validan y
qué ocurre si la validación falla.}

\subsection{Restricciones que implementan reglas de negocio}
\label{sec:restricciones}

\instruccion{Qué reglas de la \cref{sec:reglas-negocio} están sostenidas por
el propio motor de base de datos, y por qué ahí y no en el código. Es un
argumento de concurrencia: una validación en la aplicación se puede saltar si
dos peticiones llegan a la vez. Aquí va solo el mecanismo; el enunciado de la
regla vive en la sección 6.}

\ph{Restricciones e índices que sostienen las reglas y las consultas más
frecuentes.}
